\noindent\textbf{Reviewer comment:} \emph{Then, it would be helpful to add an ablation investigation demonstrating the independent contribution of the trust-aware mechanism, GRU temporal encoder, and attention-residual backbone within ART-MAPPO architecture.}

\medskip
\noindent\textbf{Response:} Thank you for this constructive recommendation. We have added a component-wise ablation in both paper scenarios. Four variants were trained separately in each scenario: complete ART-MAPPO, without trust-aware exploration, without the GRU encoder, and without the attention--residual backbone. The eight runs use the same seed, optimizer, reward, dynamics, 585-update/599,040-step budget, and scenario-specific initial-state distribution. The attention--residual control is capacity matched (732.8k versus 733.2k actor parameters), so its comparison is not explained by parameter count. Checkpoints were selected on 20 validation episodes whose seeds were disjoint from the formal test. We then froze the selected policies and ran 100 paired Monte Carlo episodes per variant and scenario (800 episodes in total), with no backpropagation or optimizer update during evaluation.

The new training curves identify complementary contributions. Removing trust causes substantially faster exploration collapse: final policy entropy changes from 1.159 to 0.517 in Case~1 and from 1.301 to 0.756 in Case~2, whereas the trust-enabled policy retains a guided-action fraction of 0.535/0.523. This directly supports the stated role of trust-aware modulation in balancing the analytical guidance prior and learned exploration. Removing the GRU increases final-window critic loss from 0.0583 to 0.3095 in Case~1 ($5.31\times$) and from 0.0348 to 0.2330 in Case~2 ($6.69\times$), together with substantially larger loss variance; hence, its clearest independent contribution is stable temporal value estimation. The capacity-matched no-attention--residual control nearly overlaps the complete model in return AUC and final critic loss under nominal conditions, but has lower final entropy in Case~2 (1.192 versus 1.301) and a slightly larger terminal normal load in Case~1 (mean 0.0340~g versus 0.0316~g). We therefore describe this backbone as providing a modest representation and safety margin under the tested nominal conditions rather than claiming a universal advantage unsupported by the data.

All eight frozen-policy cells achieve 100/100 target coverage, cooperative success, and strict mission success (Wilson 95\% interval: 96.3--100\%; exact paired McNemar $p=1$). This success-rate ceiling is informative: the analytical guidance layer and interception redundancy allow every architecture to complete the nominal tasks, so binary success alone cannot distinguish the modules. We therefore also report the four continuous metrics used in the simulation section. In Case~1, the full model shortens engagement duration by 0.419~s relative to no trust, 0.108~s relative to no GRU, and 0.032~s relative to no attention--residual (paired 95\% bootstrap intervals exclude zero). It also reduces terminal normal load relative to no GRU and no attention--residual by approximately 0.00238~g. These gains coexist with a 0.0067--0.0111~s larger group timing error for the full model, which remains far below the 0.5-s cooperative threshold. In Case~2, the continuous distributions largely overlap; the main paired operational effect is a 0.0278-s shorter duration than no GRU, whereas the other differences are small or their confidence intervals include zero.

Accordingly, the revision does not impose a forced ranking in which the full model must win every nominal metric. It provides the requested component-level evidence: trust-aware exploration preserves controlled exploration, the GRU stabilizes temporal critic learning, and the attention--residual backbone provides a capacity-controlled representation/safety margin. The common nominal success documents redundancy, while the paired continuous metrics quantify the magnitude and limits of each operational effect. We have added training-return, critic-loss, and policy-entropy comparisons, the 100-trial success summary, boxplots of $E_{co\text{-}time}$, $E_n$, $E_{miss}$, and $E_t$, and an explicit statement that the training curves are descriptive because one training seed was used per cell.
